% Generated draft: requires mathematical and source review.
\documentclass[11pt]{article}
\usepackage[margin=25mm]{geometry}
\usepackage{fontspec}
\setmainfont{Latin Modern Roman}
\usepackage{amsmath,amssymb,mathtools,unicode-math}
\setmathfont{Latin Modern Math}
\usepackage{enumitem,needspace,xurl,hyperref}
\setlength{\parindent}{0pt}
\setlength{\parskip}{5pt}
\setlength{\emergencystretch}{3em}
\newcommand{\N}{\mathbb{N}}
\newcommand{\Z}{\mathbb{Z}}
\newcommand{\Q}{\mathbb{Q}}
\newcommand{\R}{\mathbb{R}}
\newcommand{\C}{\mathbb{C}}
\newcommand{\F}{\mathbb{F}}
\newcommand{\A}{\mathbb{A}}
\newcommand{\PP}{\mathbb{P}}
\newcommand{\E}{\mathbb{E}}
\newcommand{\eps}{\varepsilon}
\newcommand{\dd}{\,\mathrm{d}}
\begin{document}
\section*{UnsolvedMath ID: 30004775}
\textbf{Source problem code:} catalogue entry
\subsection{Definitions and mathematical statement}
The supplied record attributes the following conjectural identity to Hiraga--Ichino--Ikeda. Let \(H\) denote the group intended by the source, and write \(\operatorname{Disc}(H)\) for its intended class of discrete-series representations. The record does not specify the underlying field, the hypotheses on \(H\), or whether square integrability is taken on \(H\) or modulo its centre. For \(\sigma\in\operatorname{Disc}(H)\), let \(d_H(\sigma)\) denote its formal degree, with a Haar-measure normalization that remains to be specified. The symbols \(\rho_\sigma\), \(S_\sigma\), and \(\gamma(0,\sigma,\operatorname{Ad}_H)\) are retained exactly as in the supplied statement; their constructions and conventions are not supplied. In particular, interpreting \(\dim(\rho_\sigma)\) and \(|S_\sigma|\) requires specifying the representation \(\rho_\sigma\) and the finite group \(S_\sigma\).

The requested identity is
\[
\forall\,\sigma\in\operatorname{Disc}(H),\qquad
 d_H(\sigma)
 =\frac{\dim(\rho_\sigma)}{|S_\sigma|}
 \left|\gamma\!\left(0,\sigma,\operatorname{Ad}_H\right)\right|.
\]
Here the outer absolute-value signs denote the absolute value of the gamma-factor value. The intended adjoint representation, parameter conventions, auxiliary choices for the gamma factor, and interpretation at \(s=0\) must be fixed from the source. Thus this is a faithful schematic statement of the supplied identity, not a fully specified conjecture: the missing hypotheses and normalizations prevent an unconditional precise formulation from this record alone. The background also mentions a general tempered extension, but supplies no statement of that extension; the displayed quantifier concerns only discrete-series representations.
\subsection{Short English statement}
The Hiraga--Ichino--Ikeda conjectural formula expresses the formal degree of every intended discrete-series representation of \(H\) using a dimension, a finite-group order, and the absolute value of an adjoint gamma factor at zero. The supplied record leaves the group hypotheses and essential definitions and normalizations unspecified, so this statement-only draft requires source clarification and human mathematical review.
\subsection{Sources}
\begin{itemize}
\item Local UnsolvedMath registry, supplied source record for permanent ID \(30004775\), source problem code \texttt{catalogue entry}, in the Oberwolfach Reports Open Problems collection. The record attributes the conjecture to Hiraga--Ichino--Ikeda, lists proposed year \(2021\), and supplies no separate proposer name.
\item Supplied report citation: \emph{Automorphic Forms, Geometry and Arithmetic (2022)}, Oberwolfach Reports, pp.\ 2089--2156, DOI \texttt{10.4171/owr/2021/39}. \url{https://doi.org/10.4171/owr/2021/39}. The background describes Oberwolfach Reports as documenting mathematical workshops and their open-problem discussions; the literature-source label is ``Oberwolfach Report 39/2021.'' The supplied citation's year \(2022\) and the report label's year \(2021\) are retained without resolving their relationship.
\item Supplied literature reference: \emph{Hiraga--Ichino--Ikeda conjecture for classical groups}. \url{https://arxiv.org/abs/2508.08470}. No authors or further bibliographic details for this reference are supplied.
\item Literature assessment checked at \texttt{2026-08-22}: the supplied dated curation labels the problem partially solved and states that the formal-degree conjecture is known for Archimedean groups and many \(p\)-adic cases, including general linear, unitary, odd orthogonal, regular, and unipotent families, while the general tempered extension remains unresolved. This is supplied dated curation derived from the OWR record, not an independently verified assessment; its time-sensitive status and cited claims require independent verification.
\item The supplied statement metadata carries the label \texttt{corrected\_verified} and reports that the standalone statement and its hypotheses were checked against the official Oberwolfach report. This is a supplied verification claim, not verification performed for this draft.
\end{itemize}
\end{document}
